\newpage

\begin{center}
\textbf{APÊNDICE A - EXEMPLO DE REQUISIÇÃO E RESPOSTA DA API DE ROTEAMENTO}
\end{center}

\vspace{0.5cm}

Os trechos a seguir ilustram um exemplo de requisição e de resposta do \textit{endpoint} POST /api/v1/route do \textit{DyraSQL Core}, descrito na Seção~\ref{sec:metodologia}.

\begin{center}
\textbf{Quadro 1 - Requisição do \textit{endpoint} POST /api/v1/route}
\end{center}

\begin{framed}
\begin{verbatim}
{
  "query": "SELECT ... FROM dados d JOIN tenant_info t ..."
}
\end{verbatim}
\end{framed}

\fonte{Elaborado pelo autor.}

\vspace{0.8cm}

\begin{center}
\textbf{Quadro 2 - Resposta do \textit{endpoint} POST /api/v1/route}
\end{center}

\begin{framed}
\begin{verbatim}
{
  "fingerprint": "3f2a1c...",
  "cluster": "medium",
  "score": 0.66,
  "factors": {
    "volume": 0.71,
    "complexity": 0.69,
    "historical": 0.5
  },
  "cached": false,
  "cluster_url": "http://trino-medium:8080"
}
\end{verbatim}
\end{framed}

\fonte{Elaborado pelo autor.}

\newpage

\begin{center}
\textbf{APÊNDICE B - CÓDIGO-FONTE E REPRODUÇÃO DOS EXPERIMENTOS}
\end{center}

\vspace{0.5cm}

O código-fonte do \textit{DyraSQL} está disponível em repositório público, incluindo o \textit{DyraSQL Core}, o \textit{Trino Gateway Proxy}, a orquestração em Docker Compose, os \textit{scripts} de geração e ampliação do \textit{dataset} sintético e os \textit{scripts} de medição, calibração e análise. A versão utilizada neste artigo corresponde à \textit{tag} \texttt{tcc-v1.0}, e o arquivo \texttt{REPRODUCAO.md} descreve os comandos para reproduzir os 18 casos, a calibração e a validação cruzada. O endereço do repositório é apresentado no Quadro 3.

\begin{center}
\textbf{Quadro 3 - Endereço do repositório do \textit{DyraSQL}}
\end{center}

\begin{framed}
\begin{verbatim}
https://github.com/trababalhos-puc/tcc-puc-dyrasql
tag: tcc-v1.0
\end{verbatim}
\end{framed}

\fonte{Elaborado pelo autor.}

\newpage

\begin{center}
\textbf{APÊNDICE C - TRECHO REAL DO RETORNO DO \textit{EXPLAIN (TYPE IO)}}
\end{center}

\vspace{0.5cm}

O Quadro 4 reproduz, do retorno do \textit{EXPLAIN (TYPE IO)} para a consulta do Algoritmo \ref{alg:consulta-vendas}, a entrada referente à tabela \textit{dados}, obtida no ambiente local (\textit{Trino} 484). O \textit{Score} utiliza apenas \textit{outputSizeInBytes} e \textit{outputRowCount}, somados entre as tabelas lidas. O campo \textit{cpuCost} é retornado, mas não é utilizado; o campo \textit{maxMemory} é informado como zero na leitura da tabela e como indefinido (\textit{NaN}) na estimativa do plano, razão pela qual o volume lido é empregado como indicador de demanda, e não uma estimativa de memória.

\begin{center}
\textbf{Quadro 4 - Entrada da tabela \textit{dados} no retorno do \textit{EXPLAIN (TYPE IO)}}
\end{center}

\begin{framed}
\begin{verbatim}
{
  "table": {
    "catalog": "iceberg",
    "schemaTable": {
      "schema": "analytics",
      "table": "dados"
    }
  },
  "constraint": {
    "none": false,
    "columnConstraints": [
      { "columnName": "date",
        "type": "timestamp(6)",
        "domain": {...} },
      { "columnName": "status",
        "type": "varchar",
        "domain": {...} }
    ]
  },
  "estimate": {
    "outputRowCount": 13824000.0,
    "outputSizeInBytes": 263178236.0,
    "cpuCost": 263178236.0,
    "maxMemory": 0.0,
    "networkCost": 0.0
  }
}
\end{verbatim}
\end{framed}

\fonte{Elaborado pelo autor. Os intervalos de cada filtro (\textit{domain}) foram omitidos.}
